% ============================================================
%  前置：第二版序言（p004）、第一版序言（p005–006）、符号（p007）
%  自 300dpi 页面转录；序言脚注①照原书文献注转录
% ============================================================
\chapter*{第二版序言}
\addcontentsline{toc}{chapter}{第二版序言}

在为本书准备新版时，非常遗憾的是我的合作者和朋友 Е.\,М.\,栗弗席兹已于 1985 年辞世了. 于是，我不得不独自完成此项工作. 虽然，近些年来物质凝聚态理论取得许多重要进展. 我并不认为需要从根本上改写本书. 书中讲到的理论基础及对其阐述，似乎是经住了时间的考验. 另一方面也要做些增补和修改.

增补 §35$^{*}$ 是讲超流膜的 В.\,Л.\,Березинский（1970），J.\,M.\,Kosterlitz, D.\,Thouless（1972）相变理论的. 此理论为二维系统中的过程建立了现代概念基础. 它极大丰富我们对相变本质的认识.

讲述磁性的第七章做了重大修改. 特别是增补 §74$^{*}$，其中用 T.\,Holstein, H.\,Primakoff（1940）方法研究反铁磁性. 此方法对磁性微观理论有重要意义. 因此，书中含此方法显得至关重要.

近年来，物理学领域最重要的成就之一便是发现高温超导体. 然而，实验物理学家和理论物理学家虽然加倍努力，这个有趣现象的本质仍有许多不明之处. 甚至何种相互作用造成向超导态的转变也不清楚. 在这种情况下，对此问题我只好做些简短评论.

\begin{flushright}
Л.\,П.\,皮塔耶夫斯基

2000 年 6 月
\end{flushright}

\chapter*{第一版序言}
\addcontentsline{toc}{chapter}{第一版序言}

简单说来，提供给读者的这本《理论物理学教程》第九卷是讨论物质凝聚态的量子理论的. 本卷开始于对量子液体——玻色液体和费米液体理论的详细叙述. 该理论是 Л.\,Д.\,朗道在 Л.\,Д.\,卡皮查的实验发现之后创建起来的，是当今理论物理学的一个独立分支. 这个理论的重要性，与其归因于在液态氦同位素中发生的那些有趣现象，不如说由于量子液体及其谱的概念实质上是宏观物体量子描述的基础.

例如，为了深入理解金属的性质，必须把金属中的电子看作费米液体. 但是电子液体的性质因晶格的存在而复杂化，因而对均匀且各向同性液体这种较简单的情况进行预备性研究，乃是建立理论的必需步骤. 金属的超导性也是一样，如不预先了解较简单的玻色液体超流理论，便难于清楚地把金属的超导性理解为电子液体的超流性.

格林函数方法，是现代统计物理学数学工具不可缺少的部分. 这绝不仅仅由于计算格林函数的图技术所带来的计算上的方便. 问题首先在于，格林函数能直接确定出物体的元激发谱，因此，它是最自然地描述这些激发性质的语言. 所以，对于方法问题，即宏观物体的格林函数理论，本卷相当注意. 虽然，方法的基本思想对所有的系统都一样，但在不同的情况，图技术的具体形式是有区别的. 因此，自然要以各向同性的量子液体为例来发展这些方法，这里，方法的实质以纯正的形式表露出来，没有空间非均匀性和存在几种粒子等等引起的复杂情况.

根据同样的原因，我们用简单的、具有微弱相互作用的、各向同性的费米气体模型，来阐述超导性的宏观理论，而避开由于存在库仑相互作用晶格所引起的复杂情况.

关于本书讲述晶格中的电子和磁性理论的各章节，我们再一次强调，本书是理论物理学教程的一部分，而绝无代替固体理论教程的使命. 按照这个原则，这里只研究最一般性的问题，既不要求具体实验数据的利用，也不涉及没有明确理论基础的计算方法. 我们也提醒一下，有关固体的动理学性质的内容与本卷无关，该部分拟在本教程的下一卷即最末一卷中研究.

最后，在本书中也叙述了实物介质中的电磁涨落理论和流体动力学涨落理论. 前者，原已包含在第八卷中. 现在把它移到本卷里，是出于需要运用格林函数，以使整个理论具有简单的和便于应用的形式. 此外，将电磁涨落和流体动力学涨落在同一卷里研究，也显得自然.

Л.\,Д.\,朗道不在本书的实际作者之列. 但是，读者容易看出，在本书的正文中多么频繁地遇到他的名字，因为这里所讲述的大部分成果是属于他个人的或同他的学生们合作取得的. 同他共事多年，我们有理由相信，我们做到了忠实地反映出他对这些问题的观点——当然，也顾及到从悲惨地中断他的事业那天以来的十五年间所增添的新问题.

我们诚恳地向 А.\,Ф.\,安德列耶夫（А.\,Ф.\,Андреев）、И.\,Е.\,加洛辛斯基（И.\,Е.\,Дзялошинский）和 И.\,М.\,栗弗席兹（И.\,М.\,Лифшиц）表示感谢，他们对本书中所研究的问题经常进行讨论. 我们也从著名的 А.\,А.\,阿布里科索夫（А.\,А.\,Абрикосов）、Л.\,П.\,戈里科夫（Л.\,П.\,Горьков）和 И.\,Е.\,加洛辛斯基（И.\,Е.\,Дзялошинский）的书中得到许多益处，他们的书\footnote{А.\,А.\,阿布里科索夫，Л.\,П.\,戈里科夫，И.\,Е.\,加洛辛斯基. 统计物理学中的量子场论方法. 北京：科学出版社，1963. 英译本：A.\,A.\,Abrikosov, L.\,P.\,Gorkov, I.\,E.\,Dzyaloshinski. Methods of Quantum Field Theory in Statistical Physics. Prentice-Hall, Englewood cliffs, N.\,J., 1963.}是物理学文献中讲述统计物理学新方法的首批著作之一. 最后，我们向 Л.\,П.\,戈里科夫（Л.\,П.\,Горьков）和 Ю.\,Л.\,克利蒙托维奇（Ю.\,Л.\,Климонтович）致以谢意，他们阅读了本书的原稿，并提出了许多意见.

\begin{flushright}
Е.\,М.\,栗弗席兹，Л.\,П.\,皮塔耶夫斯基

1977 年 4 月
\end{flushright}

\chapter*{符\quad 号}
\addcontentsline{toc}{chapter}{符号}

矢量的下标，用拉丁字母 $i,k,\cdots$ 标记. 自旋的下标，用希腊字母 $\alpha,\beta,\cdots$ 标记. 所有两次重复出现的下标，是指求和.

``4 维矢量''（见 46 页的注解），标以大写字母 $X,P,\cdots$.

体积元为 $\mathrm{d}V$ 或 $\mathrm{d}^3x$.

量值从上或从下趋近于零时的极限，标以 $+0$ 或 $-0$.

算符用带尖角\^{ }\ 的字母标记.

哈密顿量（哈密顿算符）为 $\hat{H}$，$\hat{H}'=\hat{H}-\mu\hat{N}$. 微扰算符是 $\hat{V}$.

薛定谔绘景中的 $\psi$ 算符为 $\hat{\psi},\hat{\psi}^{+}$；海森伯绘景中的 $\psi$ 算符是 $\hat{\Psi},\hat{\Psi}^{+}$；松原（Matsubara）绘景中的 $\psi$ 算符为 $\hat{\Psi}^{\mathrm{M}},\hat{\Psi}^{\mathrm{M}+}$.

格林函数为 $G,D$. 温度格林函数为 $\mathcal{G},\mathcal{D}$.

热力学量的标记同第五卷一样，其中有：温度为 $T$，体积是 $V$，压强是 $P$，化学势为 $\mu$.

磁场强度和磁感强度分别为 $H$ 和 $B$；外磁场是 $\mathcal{H}$.

引用本教程其它各卷的节和公式，以罗马数字按卷别表示：Ⅰ——《力学》，1988；Ⅱ——《场论》，1989；Ⅲ——《量子力学》，1989；Ⅳ——《量子电动力学》，1989；Ⅴ——《统计物理学Ⅰ》，1995；Ⅵ——《流体力学》，1986；Ⅶ——《连续介质电动力学》，1982.
